\chapter{Non-pluripolar Monge--Amp\`ere measures}
\label{chap:products}

Throughout this chapter, $\omega$ is a Hermitian form on $X$ and $\mathrm{d}V$ is a smooth positive probability volume form used for $L^1$ norms and integrals. As in \cref{sec:psh-bounded-products}, we write $\theta_v=\theta+\ddc v$, and $\mathrm{MA}_\theta(v)=(\theta+\ddc v)^n$ when $v$ is bounded; convergence in capacity refers to the capacity $\operatorname{Cap}_\omega$ of \cref{sec:hermitian-capacity}.

For an unbounded $\theta$-plurisubharmonic function $u$, the Bedford--Taylor products of \cref{prop:bounded-products} are not defined. On the set $\{u>-t\}$, however, $u$ coincides with its bounded truncation $\max(u,-t)$, and since bounded products are plurifinely local, the restricted measures are compatible as $t$ increases. Their limit is a measure on the complement of the pole set, the non-pluripolar Monge--Amp\`ere measure of $u$. For a general smooth background $\theta$, the truncation need not be $\theta$-plurisubharmonic, so positivity of the restricted products requires a local argument. The total mass of the limit may also be infinite. For a Hermitian background, finiteness of this mass for every potential is equivalent to a uniform bound on the Monge--Amp\`ere masses of bounded potentials, the bounded mass property. When the property fails, we construct a potential of infinite mass whose multiplier ideals are all trivial; its complex Hessian even has integrable coefficients, although its determinant is not integrable.

\section{Truncation and non-pluripolar measures}

Let $\theta$ be a smooth real $(1,1)$-form and $u\in\PSH(X,\theta)$. For $t>0$, put
\[
 u_t=\max(u,-t),\qquad D_t=\{u>-t\}.
\]
Choose a Hermitian form $\Omega\geq\theta$. Then $u_t$ is bounded and $\Omega$-plurisubharmonic, so \cref{prop:bounded-products} defines the product $(\theta+\ddc u_t)^n$ as a signed measure. It need not be positive on $X$, but its restriction to the plurifinely open set $D_t$ is positive. The proof replaces $u_t$ near each point by a local truncation that preserves positivity up to an arbitrarily small smooth error.

\Needspace{14\baselineskip}
\begin{lemma}\label{lem:restricted-positive-product}
With the notation above,
\[
 0\leq\mathbf1_{D_t}(\theta+\ddc u_t)^n
 \leq\mathbf1_{D_t}(\Omega+\ddc u_t)^n.
\]
More generally, let $\theta_1,\ldots,\theta_n$ be smooth real $(1,1)$-forms, let $u^i\in\PSH(X,\theta_i)$, and let $D$ be the set where $u^i>-t$ for every $i$. Then the mixed product of the currents $\theta_i+\ddc\max(u^i,-t)$ is positive on $D$.
\end{lemma}
\Needspace{7\baselineskip}
\begin{proof}
Work in a coordinate ball with a fixed positive constant form $\kappa$, and fix $\epsilon>0$. Every point has a smaller neighborhood carrying a constant real form $\theta'$ such that
\[
 -\epsilon\kappa\leq\theta-\theta'\leq\epsilon\kappa.
\]
Take a smooth local potential $r$ for $\theta'+\epsilon\kappa$. Then $u+r$ is plurisubharmonic. On a relatively compact smaller neighborhood choose $k>t+\sup|r|$ and set
\[
 h_k=\max(u+r,-k)-r.
\]
This function is locally bounded and equals $u_t=u$ on $D_t$. Moreover,
\[
 \theta+\ddc h_k+2\epsilon\kappa
 =\ddc\max(u+r,-k)+(\theta-\theta'+\epsilon\kappa)
\]
is the sum of a positive closed current with bounded potential and a smooth positive form, so its bounded products are positive. By plurifine locality, the restriction of $(\theta+\ddc u_t+2\epsilon\kappa)^n$ to $D_t$ is positive on this neighborhood.

These neighborhoods cover the coordinate ball. Now let $\epsilon\downarrow0$. The product with the fixed bounded potential $u_t$ is a polynomial in $\epsilon$ whose coefficients have locally finite variation, so it converges in local total variation, also after restriction to $D_t$. This proves the first inequality. The same construction, applied to each factor, gives the positivity of mixed products.

For the second inequality, expand the difference of the two powers by the telescoping identity, using the smooth positive form $\Omega-\theta$. Every summand is a mixed product of factors that are positive on $D_t$, as just proved. The difference is therefore nonnegative there.
\end{proof}

\Needspace{16\baselineskip}
\begin{theorem}[Extended non-pluripolar measure]\label{thm:arbitrary-extended-ma}\index{Monge--Ampere product@Monge--Amp\`ere product!extended non-pluripolar}\index{Non-pluripolar measure!extended}\label{prop:extended-measure}
Let $\theta$ be a smooth real $(1,1)$-form and $u\in\PSH(X,\theta)$. The measures $\mathbf1_{D_t}(\theta+\ddc u_t)^n$ increase with $t$. Their limit
\begin{equation}\label{eq:extended-ma}
 \mathrm{MA}^{\mathrm{ext}}_\theta(u)
 :=\lim_{t\to+\infty}\mathbf1_{\{u>-t\}}
       (\theta+\ddc\max(u,-t))^n
\end{equation}
is a positive Borel measure that puts no mass on pluripolar sets; it is finite on each $D_t$ and $\sigma$-finite on $X$. It depends only on the current $\theta+\ddc u$; in particular, it is unchanged when a constant is added to $u$. If $\Omega$ is a Hermitian form with $\theta\leq\Omega$, then
\[
 \mathrm{MA}^{\mathrm{ext}}_\theta(u)
 \leq\mathrm{MA}^{\mathrm{ext}}_\Omega(u).
\]
The total mass of $\mathrm{MA}^{\mathrm{ext}}_\theta(u)$ may be infinite.
\end{theorem}
\begin{proof}
For $s>t$, plurifine locality on $D_t\subset D_s$ gives
\[
 \mathbf1_{D_t}(\theta+\ddc u_s)^n
 =\mathbf1_{D_t}(\theta+\ddc u_t)^n.
\]
Together with \cref{lem:restricted-positive-product}, this proves that the measures increase. Their increasing limit is countably additive by monotone convergence, and it vanishes on pluripolar sets because each bounded product does. The comparison with $\Omega$ passes to the limit. By the compatibility just proved,
\[
 \mathbf1_{D_t}\mathrm{MA}^{\mathrm{ext}}_\theta(u)
 =\mathbf1_{D_t}(\theta+\ddc u_t)^n,
\]
and the mass of this measure is bounded by that of the bounded product $(\Omega+\ddc u_t)^n$. The sets $D_t$ exhaust the complement of $\{u=-\infty\}$, which has zero mass, so the limit is $\sigma$-finite.

For a smooth real function $f$, consider $\theta'=\theta+\ddc f$ and $u'=u-f$. On a common superlevel set of $u$ and $u'$, sufficiently deep truncations of $u$ and $u'$ define the same current with bounded potentials, so their restricted products agree. These sets exhaust the complement of the common pole set, which proves invariance under this change of decomposition. By the elliptic-regularity argument in the proof of \cref{prop:bounded-products}, any two smooth decompositions of the same current differ in this way, and adding a constant is a special case.
\end{proof}

We call $\mathrm{MA}^{\mathrm{ext}}_\theta(u)$ the \emph{extended non-pluripolar Monge--Amp\`ere measure} of $u$;\index{Extended non-pluripolar measure|see{Non-pluripolar measure}} the word \emph{extended} indicates that its total mass may be infinite. On compact $X$ it defines a current of order zero exactly when its total mass is finite. In that case it is the non-pluripolar Monge--Amp\`ere measure in the usual sense, namely the extension by zero across $\{u=-\infty\}$ of the plurifinely local bounded products. If $u$ is bounded, then $D_t=X$ for large $t$, and $\mathrm{MA}^{\mathrm{ext}}_\theta(u)=\mathrm{MA}_\theta(u)$. In accordance with the usual notation for non-pluripolar products, we also write $\langle(\theta+\ddc u)^n\rangle$ for $\mathrm{MA}^{\mathrm{ext}}_\theta(u)$; lower powers $\langle(\theta+\ddc u)^k\rangle$ are defined by the same truncation where they are used.\index{Non-pluripolar measure!angle bracket notation}

\section{Contact sets}

The bounded contact inequality of \cref{lem:bounded-ordered-contact} extends to unbounded potentials on common finite levels. As in the previous section, local positive corrections are needed, because the global truncations need not remain $\theta$-plurisubharmonic.

\Needspace{12\baselineskip}
\begin{proposition}\label{prop:ordered-contact-general}\index{Contact set!ordered potentials}
Let $\theta$ be a smooth real $(1,1)$-form and let $u,v\in\PSH(X,\theta)$ satisfy $u\leq v$. Then
\[
 \mathbf1_{\{u=v\}}\mathrm{MA}^{\mathrm{ext}}_\theta(u)
 \leq\mathbf1_{\{u=v\}}\mathrm{MA}^{\mathrm{ext}}_\theta(v).
\]
\end{proposition}
\begin{proof}
Fix $t>0$, put $D_t=\{u>-t\}$, and choose a Hermitian form $\Omega\geq\theta$. The truncations $u_t=\max(u,-t)$ and $v_t=\max(v,-t)$ are bounded and $\Omega$-plurisubharmonic, so their signed products with background $\theta$ are defined. Fix $\epsilon>0$ and a point of a coordinate ball, and let $\kappa$, $r$ and $k>t+\sup|r|$ be as in the proof of \cref{lem:restricted-positive-product}. On a neighborhood of the point, the bounded functions
\[
 a_k=\max(u+r,-k)-r,\qquad
 b_k=\max(v+r,-k)-r
\]
satisfy $a_k\leq b_k$, and as shown there, $\theta+2\epsilon\kappa+\ddc a_k\geq0$ and $\theta+2\epsilon\kappa+\ddc b_k\geq0$. On $D_t$ we have $v\geq u>-t$, so $a_k=u_t$ and $b_k=v_t$ there. The top-degree case of \cref{lem:bounded-ordered-contact}, with background $\theta+2\epsilon\kappa$, and plurifine locality give
\[
 \mathbf1_{\{u=v\}\cap D_t}(\theta+2\epsilon\kappa+\ddc u_t)^n
 \leq\mathbf1_{\{u=v\}\cap D_t}(\theta+2\epsilon\kappa+\ddc v_t)^n
\]
on this neighborhood, hence on the whole ball. As in the proof of \cref{lem:restricted-positive-product}, both sides converge in local total variation as $\epsilon\downarrow0$, which gives the same inequality for the background $\theta$. Since $D_t\subset\{v>-t\}$, \cref{thm:arbitrary-extended-ma} identifies its two sides with the restrictions of $\mathrm{MA}^{\mathrm{ext}}_\theta(u)$ and $\mathrm{MA}^{\mathrm{ext}}_\theta(v)$ to $\{u=v\}\cap D_t$. Finally let $t\to+\infty$: the sets $\{u=v\}\cap D_t$ exhaust the contact set off the pluripolar set $\{u=-\infty\}$, which has zero extended mass.
\end{proof}

\section{Weighted lower semicontinuity}

Products of uniformly bounded potentials are continuous under convergence in capacity. Multiplying the extended measure by a cutoff in the potential that vanishes below a fixed level reduces it to such bounded products, and letting the level tend to $-\infty$ gives lower semicontinuity of the extended measures, even when their masses are infinite. The weights are allowed to vary with the potentials.

\Needspace{12\baselineskip}
\begin{lemma}\label{lem:weighted-extended-lsc}\index{Non-pluripolar measure!lower semicontinuity}\label{lem:big-weighted-lsc}
Let $\theta$ be a smooth real $(1,1)$-form. Suppose $u_j,u\in\PSH(X,\theta)$ have a common upper bound and $u_j\to u$ in capacity. If $b_j,b$ are nonnegative uniformly bounded quasi-continuous functions and $b_j\to b$ in capacity, then
\[
 \int_Xb\,\mathrm{MA}^{\mathrm{ext}}_\theta(u)
 \leq\liminf_j\int_Xb_j\,\mathrm{MA}^{\mathrm{ext}}_\theta(u_j).
\]
Both sides may be infinite.
\end{lemma}
\begin{proof}
Subtracting a common constant, we may assume that all the potentials are nonpositive. Choose a Hermitian form $\Omega\geq\theta$. For $T>0$, let $g_T\colon [-\infty,0]\to[0,1]$ be a continuous nondecreasing cutoff, equal to zero on $[-\infty,-2T]$ and to one on $[-T,0]$, chosen so that $g_T$ increases pointwise to one on $(-\infty,0]$ as $T\to\infty$. Put $v_j=\max(u_j,-2T)$ and $v=\max(u,-2T)$. These are uniformly bounded $\Omega$-plurisubharmonic functions, and they converge in capacity.

Fix $T$. The weights $b_jg_T(v_j)$ are uniformly bounded and quasi-continuous, and they converge in capacity to $bg_T(v)$. The weighted capacity continuity of bounded products (\cref{sec:psh-local-capacity,sec:hermitian-capacity}), applied term by term to the local expansion of the signed products $(\theta+\ddc v_j)^n$, gives the convergence below. Since the cutoff vanishes at the truncation level $-2T$, locality gives the identity
\[
 \begin{split}
 \int_Xb_jg_T(u_j)\,\mathrm{MA}^{\mathrm{ext}}_\theta(u_j)
 &=\int_Xb_jg_T(v_j)(\theta+\ddc v_j)^n\\
 &\longrightarrow
 \int_Xbg_T(u)\,\mathrm{MA}^{\mathrm{ext}}_\theta(u).
 \end{split}
\]
The first integral is nonnegative and at most the corresponding integral without the cutoff, although the bounded product $(\theta+\ddc v_j)^n$ may be signed. Hence the limit is bounded by the lower limit in the assertion. Letting $T\to\infty$ and using monotone convergence completes the proof.
\end{proof}

\section{Density and trace inequalities}
\label{sec:extended-density-inequalities}

The mixed density inequality of \cref{lem:nonclosed-mixed-density} passes to unbounded potentials on common finite levels. It gives a lower bound for the extended measure of a convex combination of potentials and, by comparison with a smooth metric, a lower bound for the full trace measure.

\Needspace{16\baselineskip}
\begin{lemma}\label{lem:singular-convex-density}
Let $\theta$ be a smooth semipositive $(1,1)$-form and let $f,g\geq0$ be measurable functions on $X$. Suppose $u,v\in\PSH(X,\theta)$ satisfy
\[
 \mathrm{MA}^{\mathrm{ext}}_\theta(u)\geq f\,\mathrm{d}V,
 \qquad \mathrm{MA}^{\mathrm{ext}}_\theta(v)\geq g\,\mathrm{d}V.
\]
Then, for $0\leq t\leq1$,
\begin{equation}\label{eq:singular-convex-density}
 \mathrm{MA}^{\mathrm{ext}}_\theta(tu+(1-t)v)
 \geq\bigl(tf^{1/n}+(1-t)g^{1/n}\bigr)^n\,\mathrm{d}V.
\end{equation}
\end{lemma}
\begin{proof}
The cases $t=0$ and $t=1$ are the hypotheses. Let $0<t<1$, translate $u,v$ to be nonpositive, and put $u_k=\max(u,-k)$ and $v_k=\max(v,-k)$. Since $\theta$ is semipositive, these are bounded $\theta$-plurisubharmonic functions, and by locality
\[
 \mathrm{MA}_\theta(u_k)\geq\mathbf1_{\{u>-k\}}f\,\mathrm{d}V,
 \qquad
 \mathrm{MA}_\theta(v_k)\geq\mathbf1_{\{v>-k\}}g\,\mathrm{d}V.
\]
On $E_k=\{u>-k\}\cap\{v>-k\}$, the inequality \eqref{eq:interior-minkowski-background} of \cref{lem:nonclosed-mixed-density}, with $a=t$ and $b=1-t$, bounds $\mathrm{MA}_\theta(tu_k+(1-t)v_k)$ below by the measure on the right of \eqref{eq:singular-convex-density}. On the same set, plurifine locality identifies the left-hand side with the extended measure of $tu+(1-t)v$. As $k\to\infty$, the sets $E_k$ exhaust the complement of the two pole sets, which have zero smooth volume. Monotone convergence on Borel sets proves the inequality.
\end{proof}

The same mixed inequality gives a lower bound for the full trace measure of an unbounded potential. Unlike the extended measure, the full current $\omega+\ddc u$ keeps its mass on the pole set; the bound below concerns only its absolutely continuous part.

\begin{lemma}\label{lem:unbounded-gauduchon-trace-density}
Let $\gamma$ be a Hermitian form. If $u\in\PSH(X,\omega)$ and a measurable function $g\geq0$ satisfy $\mathrm{MA}^{\mathrm{ext}}_\omega(u)\geq g\,\mathrm{d}V$, then
\[
 (\omega+\ddc u)\wedge\gamma^{n-1}
 \geq(g\,\mathrm{d}V/\gamma^n)^{1/n}\gamma^n.
\]
\end{lemma}
\begin{proof}
For bounded $u$, \cref{lem:nonclosed-mixed-density}, applied in local coordinates with $T_1=\omega+\ddc u$, $T_2=\gamma$ and $k=1$, gives the asserted inequality.

In general, subtract a constant so that $u\leq0$, and put $u_k=\max(u,-k)$. Positivity and locality give
\[
 \mathrm{MA}_\omega(u_k)\geq\mathbf1_{\{u>-k\}}g\,\mathrm{d}V,
\]
so the bounded case applies to $u_k$ with the density $\mathbf1_{\{u>-k\}}g$. Since $u_k\to u$ in $L^1$, the positive measures $(\omega+\ddc u_k)\wedge\gamma^{n-1}$ converge weakly to $(\omega+\ddc u)\wedge\gamma^{n-1}$. The densities on the right increase to $(g\,\mathrm{d}V/\gamma^n)^{1/n}$ off the pole set, which has zero smooth volume. Testing against nonnegative continuous functions and using monotone convergence on the right proves the assertion.
\end{proof}

\section{The bounded mass property}

When $\omega$ is a K\"ahler form, all bounded $\omega$-plurisubharmonic functions have the same Monge--Amp\`ere mass $\int_X\omega^n$. For a general Hermitian form, the mass can depend on the potential, as in \cref{ex:torus-mass}, and the following two numbers measure its range.

\begin{definition}\label{def:bmp}\index{BMP|see{Bounded mass property}}\index{Bounded mass property}\index{Positive lower mass!Hermitian backgrounds}\index{Lower mass}\index{Upper mass}
Let $\omega$ be a Hermitian form on $X$. The \emph{lower mass} and the \emph{upper mass} of $\omega$ are
\[
 \begin{split}
 v_-(\omega)&=
 \inf_{v\in\PSH(X,\omega)\cap L^\infty(X)}
       \int_X(\omega+\ddc v)^n,\\
 v_+(\omega)&=
 \sup_{v\in\PSH(X,\omega)\cap L^\infty(X)}
       \int_X(\omega+\ddc v)^n.
 \end{split}
\]
The manifold $X$ has the \emph{bounded mass property}, abbreviated BMP, if $v_+(\omega)<+\infty$. The form $\omega$ has \emph{positive lower mass} if $v_-(\omega)>0$.
\end{definition}

Testing with $v=0$ gives $v_-(\omega)\leq\int_X\omega^n\leq v_+(\omega)$. By the next proposition, the bounded mass property does not depend on the choice of $\omega$. Since the extended measure \eqref{eq:extended-ma} is built from products of bounded potentials, the upper mass also controls the extended measures of unbounded potentials (\cref{prop:bmp-finite}).

\begin{proposition}\label{prop:bmp-metric}\index{Bounded mass property!metric independence}
The bounded mass property is independent of the Hermitian form. Moreover, $v_+(\omega)=\Vsm(\omega)$, the supremum of the masses of smooth strictly $\omega$-plurisubharmonic potentials in \cref{prop:smooth-bmp-metric}; in particular, the supremum in \cref{def:bmp} is unchanged if only smooth potentials are used.
\end{proposition}
\begin{proof}
Smooth strictly $\omega$-plurisubharmonic potentials are bounded, so $\Vsm(\omega)\leq v_+(\omega)$. Conversely, approximate a bounded potential by \cref{thm:bounded-exact-regularization}. The smooth approximants are strictly $\omega$-plurisubharmonic, decrease and remain uniformly bounded, so the continuity of bounded products along decreasing sequences gives convergence of their total masses. This proves $v_+(\omega)=\Vsm(\omega)$, and metric independence follows from \cref{prop:smooth-bmp-metric}.
\end{proof}

\begin{proposition}\label{prop:bmp-finite}
If $\theta$ is a smooth real $(1,1)$-form with $\theta\leq C\omega$ for some $C>0$, and $u\in\PSH(X,\theta)$, then
\[
 \mathrm{MA}^{\mathrm{ext}}_\theta(u)(X)\leq C^n v_+(\omega).
\]
\end{proposition}
\begin{proof}
The comparison in \cref{thm:arbitrary-extended-ma}, applied with the Hermitian form $C\omega$, gives
\[
 \mathbf1_{D_t}\mathrm{MA}^{\mathrm{ext}}_\theta(u)
 \leq\mathbf1_{D_t}(C\omega+\ddc u_t)^n.
\]
Since $u_t/C$ is bounded and $\omega$-plurisubharmonic, the right-hand side has mass at most $C^n v_+(\omega)$. Let $t\to\infty$.
\end{proof}

\Needspace{20\baselineskip}
\section{Infinite mass with trivial multiplier ideals}

If the bounded mass property fails, there are smooth $\omega$-plurisubharmonic functions $u_j$ whose masses $M_j$ tend to infinity. In a convex combination $U=\sum_ja_ju_j$, positivity of mixed products shows that each summand contributes at least $a_j^nM_j$ to the extended measure of $U$. Choosing the masses $M_j$ to outgrow $a_j^{-n}$ therefore forces infinite mass. On the other hand, each tail $\sum_{j>N}a_ju_j$ is a small multiple of a convex combination of normalized potentials, and the uniform exponential integrability of \cref{lem:uniform-exponential} applies to it at every exponent.

\begin{theorem}\label{thm:infinite-mass-witness}
If $v_+(\omega)=+\infty$, there is $U\in\PSH(X,\omega)$, $U\leq0$, such that
\[
 \mathrm{MA}^{\mathrm{ext}}_\omega(U)(X)=+\infty,
 \qquad
 \int_X\mathrm{e}^{-pU}\,\mathrm{d}V<+\infty\quad\hbox{for every }p>0.
\]
In particular, $\mathcal I(pU)=\mathcal O_X$ for every $p>0$.
\end{theorem}
\begin{proof}
Set $a_j=2^{-j}$ for $j\geq1$. By \cref{prop:bmp-metric}, we can choose smooth $\omega$-plurisubharmonic functions $u_j$ with $\sup_Xu_j=0$ and
\[
 M_j:=\int_X\omega_{u_j}^n\geq2^{j(n+1)}.
\]
The partial sums $U_N=\sum_{j\leq N}a_ju_j$ decrease and satisfy
\[
 \omega+\ddc U_N
 =\sum_{j\leq N}a_j\omega_{u_j}
     +\left(1-\sum_{j\leq N}a_j\right)\omega\geq0.
\]
Their $L^1$ norms are bounded by \eqref{eq:normalized-lone}, so their limit $U$ is $\omega$-plurisubharmonic and nonpositive.

We first prove that the mass is infinite. Fix $N$, and put
\[
 S_N=\sum_{j\leq N}a_j\omega_{u_j},\qquad
 \delta_N=\sum_{j>N}a_j=2^{-N},\qquad
 V_N=U-U_N.
\]
Passing to the limit in the remaining partial sums gives $\delta_N\omega+\ddc V_N\geq0$. For $t>0$, the function $W_{N,t}=U_N+\max(V_N,-t)$ is bounded, and its current is the sum of two positive currents with bounded potentials:
\[
 \omega+\ddc W_{N,t}
 =S_N+\bigl(\delta_N\omega+\ddc\max(V_N,-t)\bigr).
\]
On $D_t=\{U>-t\}$, the inequality $U_N\leq0$ gives $V_N\geq U>-t$, and hence $W_{N,t}=U=\max(U,-t)$ there. Plurifine locality and positivity of bounded mixed products yield
\[
 \mathbf1_{D_t}\omega_{\max(U,-t)}^n
 \geq\mathbf1_{D_t}S_N^n.
\]
The sets $D_t$ exhaust $X\setminus\{U=-\infty\}$. This pole set is pluripolar, hence has zero smooth volume, so the smooth measure $S_N^n$ gives it no mass. Letting $t\to+\infty$, we obtain
\begin{equation}\label{eq:infinite-mass-partial-sums}
 \mathrm{MA}^{\mathrm{ext}}_\omega(U)\geq S_N^n.
\end{equation}
For $j\leq N$, positivity gives $S_N^n\geq a_j^n\omega_{u_j}^n$, so the mass of $U$ is at least $a_j^nM_j\geq2^j$. Since $j$ is arbitrary, the mass is infinite.

For the integrability, fix $p>0$, let $\alpha$ be the exponent of \cref{lem:uniform-exponential}, and take $N$ so large that
$p\delta\leq\alpha$, where $\delta=\sum_{j>N}a_j=2^{-N}$. Write
$U=A_N+V_N$ with $A_N=\sum_{j\leq N}a_ju_j$ smooth.
For a finite tail ending at $m>N$, Jensen's inequality for the convex exponential function, with the unused weight assigned to the value zero, gives
\begin{align*}
 \exp\!\left(-p\sum_{N<j\leq m}a_ju_j\right)
 &\leq\sum_{N<j\leq m}\frac{a_j}{\delta}
                  \mathrm{e}^{-p\delta u_j}
       +1-\sum_{N<j\leq m}\frac{a_j}{\delta}.
\end{align*}
Because $u_j\leq0$ and $p\delta\leq\alpha$, the integrals of the exponentials on the right are uniformly bounded. Integrating and letting $m\to+\infty$ by monotone convergence, we obtain $\int_X\mathrm{e}^{-pV_N}\mathrm{d}V<+\infty$. Multiplication by the bounded smooth function $\mathrm{e}^{-pA_N}$ proves $\int_X\mathrm{e}^{-pU}\mathrm{d}V<+\infty$.

Consequently, the constant germ $1$ is locally integrable against each weight $\mathrm{e}^{-pU}$, and every multiplier ideal $\mathcal I(pU)$ is the whole structure sheaf.
\end{proof}

\Needspace{11\baselineskip}
\begin{corollary}\label{cor:bmp-characterization}
\index{Bounded mass property!finiteness of all potential masses}
The bounded mass property holds if and only if
$\mathrm{MA}^{\mathrm{ext}}_\omega(u)(X)<+\infty$ for every $u\in\PSH(X,\omega)$.
If the bounded mass property fails, there are potentials with infinite extended mass and trivial multiplier ideals $\mathcal I(pu)=\mathcal O_X$, $p>0$, arbitrarily close to zero in $L^1$.
\end{corollary}
\begin{proof}
The first assertion follows from \cref{prop:bmp-finite,thm:infinite-mass-witness}. For the second, take $\varepsilon U$, $0<\varepsilon\leq1$, with $U$ as in \cref{thm:infinite-mass-witness}. It is $\omega$-plurisubharmonic and has the required exponential integrability. On $D_t=\{U>-t\}$, its truncation at level $\varepsilon t$ is $\varepsilon\max(U,-t)$, and
\[
 \omega+\ddc\bigl(\varepsilon\max(U,-t)\bigr)
 =\varepsilon\omega_{\max(U,-t)}+(1-\varepsilon)\omega.
\]
Positivity of mixed products bounds the restricted measure of $\varepsilon U$ below by $\varepsilon^n$ times that of $U$. Its limiting mass is therefore infinite. Finally, $\|\varepsilon U\|_{L^1}\to0$ as $\varepsilon\to0$.
\end{proof}

\section{Integrable Hessians and infinite mass}
\label{sec:integrable-infinite-mass}

The construction in \cref{thm:infinite-mass-witness} has a stronger linear convergence property. Its limiting current has integrable coefficients, and the smooth approximating currents converge in the coefficient $L^1$ norm. Nevertheless their determinant masses tend to infinity. Thus the obstruction is already present for an ordinary measurable matrix of integrable coefficients.

Let $X$ be a connected compact complex manifold of dimension $n$, and fix a Hermitian form $\omega$ and a Gauduchon form $\gamma$ on $X$. Let $v_+(\omega)$ denote the upper bounded-potential mass of \cref{def:bmp}, and use $\mathrm{MA}^{\mathrm{ext}}_\omega$ for the extended measure of \cref{thm:arbitrary-extended-ma}. The notation $\mathcal I(pu)$ denotes the multiplier ideal of the weight $pu$, and $\mathcal O_X$ is the structure sheaf. Put
\[
 M=\int_X\omega\wedge\gamma^{n-1}.
\]
For a real $(1,1)$-form $T$ with measurable coefficients, let $A_T$ be its Hermitian matrix relative to $\gamma$, and let $|A_T|=(A_T^2)^{1/2}$ be its positive matrix absolute value. Use the norm
\[
 \|T\|_{L^1(\gamma)}
 =\frac1n\int_X\operatorname{tr}|A_T|\,\gamma^n.
\]
For $T\geq0$ this equals $\int_XT\wedge\gamma^{n-1}$. Any smooth background norm gives the same space of integrable coefficients on compact $X$. We write $\det_\gamma T=\det A_T$, defined almost everywhere when the coefficients are integrable.

\Needspace{24\baselineskip}
\begin{proposition}\label{prop:integrable-infinite-mass}
Suppose $v_+(\omega)=+\infty$. The potential $U$ and smooth decreasing approximants $U_N$ in the proof of \cref{thm:infinite-mass-witness} satisfy the following additional conclusions. The positive current $T=\omega+\ddc U$ has integrable coefficients, and
\[
 \|\omega+\ddc U_N-T\|_{L^1(\gamma)}
 \leq2^{1-N}M.
\]
Its pointwise determinant satisfies
\[
 \mathrm{MA}^{\mathrm{ext}}_\omega(U)
 \geq (\det_\gamma T)\gamma^n,
 \qquad
 \int_X(\det_\gamma T)\gamma^n=+\infty.
\]
Moreover,
\[
 \int_X(\omega+\ddc U_N)^n\longrightarrow+\infty.
\]
The determinant measure in the statement is allowed infinite mass; equality with the non-pluripolar measure is not required.
\end{proposition}
\begin{proof}
Use the notation of that construction:
\[
 a_j=2^{-j},\quad T_j=\omega+\ddc u_j\geq0,\quad
 U_N=\sum_{j\leq N}a_ju_j\downarrow U,\quad
 \int_XT_j^n\geq2^{j(n+1)}.
\]
Every $u_j$ is smooth and nonpositive. Put
\[
 S_N=\sum_{j\leq N}a_jT_j,\qquad
 \delta_N=\sum_{j>N}a_j=2^{-N}.
\]
The smooth positive matrices $S_N$ increase pointwise. The Gauduchon identity gives
\[
 \int_XS_N\wedge\gamma^{n-1}=(1-\delta_N)M.
\]
Monotone convergence of their traces shows that the matrix series has a finite positive limit $S$ almost everywhere. Its coefficients are integrable: in a $\gamma$-unitary frame the absolute value of every coefficient is bounded by the trace. Positivity and the same trace identity give
\[
 \|S-S_N\|_{L^1(\gamma)}=\delta_N M.
\]
Thus $S_N\to S$ also in distributions. Since
\[
 \omega+\ddc U_N=S_N+\delta_N\omega
 \longrightarrow\omega+\ddc U
\]
in distributions by $L^1$ convergence of the potentials, we obtain $T=S$. The triangle inequality now gives
\[
 \|\omega+\ddc U_N-T\|_{L^1(\gamma)}
 \leq\|S_N-S\|_{L^1(\gamma)}
       +\delta_N\|\omega\|_{L^1(\gamma)}
 =2\delta_N M.
\]

We next compare the non-pluripolar measure with the pointwise determinant, using only bounded products. Fix $N$ and set $V_N=U-U_N$. The remaining partial sums show that
\[
 \delta_N\omega+\ddc V_N\geq0,\qquad V_N\leq0.
\]
For $t>0$, define
\[
 W_{N,t}=U_N+\max(V_N,-t).
\]
It is bounded, and its current is the sum of two positive factors:
\[
 \omega+\ddc W_{N,t}
 =S_N+\bigl(\delta_N\omega+\ddc\max(V_N,-t)\bigr).
\]
On $D_t=\{U>-t\}$, the inequality $U_N\leq0$ gives $V_N\geq U>-t$. Consequently $W_{N,t}=U=\max(U,-t)$ there. Locality and positivity of bounded products imply
\[
 \mathbf1_{D_t}\mathrm{MA}^{\mathrm{ext}}_\omega(U)
 \geq\mathbf1_{D_t}S_N^n.
\]
The smooth measure $S_N^n$ ignores the pluripolar pole set of $U$. Letting $t\to\infty$ proves
$\mathrm{MA}^{\mathrm{ext}}_\omega(U)\geq S_N^n$.
The positive matrices $S_N$ increase to $T$ almost everywhere, so their determinants increase to $\det_\gamma T$. Monotone convergence on every Borel set gives the asserted determinant-measure inequality. This does not presume that either limiting measure is locally finite.

For every $N\geq j$, positivity gives $S_N\geq a_jT_j$. Hence both the measurable limit and the smooth approximating current satisfy
\[
 \int_X(\det_\gamma T)\gamma^n
 \geq a_j^n\int_XT_j^n\geq2^j,
 \qquad
 \int_X(\omega+\ddc U_N)^n
 \geq a_j^n\int_XT_j^n\geq2^j.
\]
The first inequality holds for every $j$, so the integral is infinite. For the second, every $N\geq j$ obeys the bound; this proves divergence of the full sequence, without claiming that its masses are monotone.
\end{proof}

\Needspace{28\baselineskip}
\begin{corollary}\label{cor:strong-linear-infinite-mass}
If bounded mass fails, there are nonpositive potentials $w_\varepsilon\in\PSH(X,\omega)$, $0<\varepsilon<1$, such that
\[
 \omega+\ddc w_\varepsilon\geq(1-\varepsilon)\omega,
 \qquad
 \mathcal I(pw_\varepsilon)=\mathcal O_X\quad(p>0),
\]
and both the pointwise determinant integral and the extended Monge--Amp\`ere mass are infinite. As $\varepsilon\downarrow0$, the potentials converge to zero in $L^1$ and their currents converge to $\omega$ in the coefficient $L^1$ norm. Their full Gauduchon trace measures converge to $\omega\wedge\gamma^{n-1}$ in total variation.
\end{corollary}
\begin{proof}
Take $w_\varepsilon=\varepsilon U$. Its current is
\[
 T_\varepsilon=(1-\varepsilon)\omega+\varepsilon T.
\]
The lower bound and the two $L^1$ convergence assertions follow immediately. The full trace measures have integrable densities, and their total-variation distance is at most
\[
 \varepsilon\|T-\omega\|_{L^1(\gamma)}.
\]
Every multiplier ideal is trivial by the exponential integrability already proved for $U$. Since $T_\varepsilon\geq\varepsilon T$ almost everywhere, its determinant integral is at least $\varepsilon^n$ times that of $T$, hence infinite. For the extended measure, constant truncation and mixed positivity give
\[
 \mathrm{MA}^{\mathrm{ext}}_\omega(\varepsilon U)
 \geq\varepsilon^n\mathrm{MA}^{\mathrm{ext}}_\omega(U),
\]
as in \cref{cor:bmp-characterization}. Its mass is therefore infinite as well.
\end{proof}

In this construction the full current, and therefore every smooth trace of it, is absolutely continuous with respect to smooth volume. The failure of top-degree integrability does not arise from mass of that current on a pluripolar set. Linear trace control, even with total-variation convergence and trivial multiplier ideals, must still be supplemented by nonlinear estimates.

